% previous work have shown success in 2d .. but what about 3D? 3D has always been difficult for ... reasons
%--> existing 3d approaches are either low quality or just not interactable
%--> we are curious, can we leverage the commonsense priors and spatial knowledge inherent in LLMs to generate diverse, realistic, nad interactable scenes?
%-->traditionally, indoor scene generation has been broken down into different tasks and there are specific models for each tasks 
%--> we aim to tackle the end to end process of indoor scene generation by using LLM.
% Previous work has largely compartmentalized these aspects, offering specialized models for tasks such as floor plan generation~\cite{} or object rearrangement~\cite{}. However, these task-specific models cannot generate diverse 3D scenes based on text prompts or are often limited by the constraints of a ``training set''. 
%The landscape of past work in 3D scene generation can be broadly categorized into three lines of work.
% Despite the success of generative models in generating 3D \textsc{objects}~\cite{poole2022dreamfusion, lin2023magic3d, huang2023aladdin}, the generation of 3D \textsc{environments} from text remains limited due to the scarce availability of 3D scene data and complex compositionality of a scene layout. 

% talk about recent work and their limitations: mesh-based, not interactable, require training data. not generalizable
% only focus on residential scenes
% procthor, randomly sampled, cannot control the generation, not realistic

% Generative models have recently achieved remarkable success in visual domains such as 2D images \cite{rombach2022high, ramesh2022hierarchical, saharia2022photorealistic} and 3D objects \cite{poole2022dreamfusion, lin2023magic3d, huang2023aladdin}.
% However, the landscape for generating interactive 3D environments from text is markedly less mature despite its wide applicability in AR/VR, robotics, games, films, etc.

%The creation of 3D interactive environments has broad implications for Embodied AI.
%3D environments are vital for providing realistic, diverse, and complex scenarios essential for training robust embodied agents, where the predominant approach involves training agents in simulators and then transferring onto real robots \cite{phone2proc, Kadian2019Sim2RealPD}. 

%The predominant approach in training embodied agents involves training agents in simulators~\cite{ai2thor, xia2018gibson, puig2018virtualhome, savva2019habitat} and then transferring onto real robots \cite{phone2proc, Kadian2019Sim2RealPD}. Having realistic, diverse, and complex 3D environments plays a crucial role in the success of this sim-to-real process.
The predominant approach in training embodied agents involves learning in simulators~\cite{ai2thor, xia2018gibson, puig2018virtualhome, savva2019habitat, phone2proc, Kadian2019Sim2RealPD}. Generating realistic, diverse, and interactive 3D environments plays a crucial role in the success of this process.
% Challenges in generating 3D environments are manifold: high-dimensional 3D data in limited availability, semantic consistency between different elements within a scene, complex compositionality of a scene layout, and using assets supported by an interactive simulator.

Existing Embodied AI environments are typically crafted through manual design \cite{ai2thor, deitke2020robothor,  li2023behavior, gan2020threedworld}, 3D scanning \cite{savva2019habitat, ramakrishnan2021hm3d, phone2proc}, or procedurally generated with hard-coded rules \cite{procthor}. However, these methods require considerable human effort that involves designing a complex layout, using assets supported by an interactive simulator, and placing them into scenes while ensuring semantic consistency between the different scene elements. Therefore, prior work on producing 3D environments mainly focuses on limited environment types.
To move beyond these limitations, recent works adapt 2D foundational models to generate 3D scenes from text~\cite{zhang2023text2nerf,hollein2023text2room,fridman2023scenescape}. However, these models often produce scenes with significant artifacts, such as mesh distortions, and lack the interactivity necessary for Embodied AI.
Moreover, there are models tailored for specific tasks like floor plan generation~\cite{hu2020graph2plan,shabani2023housediffusion} or object arrangement~\cite{wei2023lego,paschalidou2021atiss}. Although effective in their respective domains, they lack overall scene consistency and rely heavily on task-specific datasets.


In light of these challenges, we present \textbf{\holodeck}, a language-guided system built upon AI2-THOR \cite{ai2thor}, to automatically generate diverse, customized, and interactive 3D embodied environments from textual descriptions. Shown in Figure \ref{fig: pipeline}, given a description (e.g., \emph{a 1b1b apartment of a researcher who has a cat}), \holodeck uses a Large Language Model (GPT-4 \cite{openai2023gpt}) to design the floor plan, assign suitable materials, install the doorways and windows and arrange 3D assets coherently in the scene using constraint-based optimization. \holodeck chooses from over 50K diverse and high-quality 3D assets from Objaverse~\cite{deitke2023objaverse} to satisfy a myriad of environment descriptions.
 
%\holodeck is modularized, with each module handling different generation steps by taking responses from LLM as inputs to specify the scenes.

Motivated by the emergent abilities of Large Language Models (LLMs) \cite{wei2022emergent}, \holodeck exploits the commonsense priors and spatial knowledge inherently present in LLMs. This is exemplified in Figure \ref{fig: demo examples}, where \holodeck creates diverse scene types such as \textit{arcade}, \textit{spa} and \textit{museum}, interprets specific and abstract prompts by placing relevant objects appropriately into the scene, e.g., an ``R2-D2''\footnote{A fictional robot character in the Star Wars.} on the desk for ``a fan of Star Wars". Beyond object selection and layout design, \holodeck showcases its versatility in style customization, such as creating a scene in a ``Victorian-style" by applying appropriate textures and designs to the scene and its objects. Moreover, \holodeck demonstrates its proficiency in spatial reasoning, like devising floor plans for ``three professors' offices connected by a long hallway'' and having regular arrangements of objects in the scenes.
% \ani{Add the disability example to figure 1 and then describe it here.}
Overall, \holodeck offers a broad coverage approach to 3D environment generation, where textual prompts unlock new levels of control and flexibility in scene creation.

The effectiveness of \holodeck is assessed through its scene generation quality and applicability to Embodied AI. Through large-scale user studies involving 680 participants, we demonstrate that \holodeck significantly surpasses existing procedural baseline \procthor \cite{procthor} in generating residential scenes and achieves high-quality outputs for various scene types. For the Embodied AI experiments, we focus on \holodeck's application in aiding zero-shot object navigation in previously unseen scene types. % Our results indicate that, with the intended scene type for agent deployment as the only known information , 
We show that agents trained on scenes generated by \holodeck can navigate better in novel environments (e.g., \textit{Daycare} and \textit{Gym}) designed by experts.

%While for Embodied AI applications, when tested on the challenging object navigation tasks in novel environments (\textit{Daycare}, \textit{Gym}, etc.) designed by experts, agents fine-tuned on \holodeck-generated scenes show improved generalization capabilities.

To summarize, our contributions are three-fold: (1) We propose \holodeck, a language-guided system capable of generating diverse, customized, and interactive 3D environments based on textual descriptions; (2) The human evaluation validates \holodeck's capability of generating residential and diverse scenes with accurate asset selection and realistic layout design; (3) Our experiments demonstrate that \holodeck can aid Embodied AI agents in adapting to new scene types and objects during object navigation tasks.
%illustrate that \holodeck effectively narrows the generalization gap under scene type shifts.

% (2) Our novel object placement module utilizes spatial relational constraints combined with an optimization process (we have ablation studies to support this)

% \holodeck leverages the commonsense and spatial knowledge from LLM to generate the scene procedurally including  
% The system solves semantic consistency by sampling commonsense attributes for scene assets and commonsense relations between assets in the scene, followed by an optimization process to solve for the complex compositionality intrinsic to generating 3D scenes. 
% These mechanisms allow \holodeck to generate indoor scenes that are semantically consistent, holistically, and spatially coherent. 

% As shown in Figure \ref{fig: pipeline}, multiple specialized modules are required to generate a comprehensive, interactable indoor 3D scene. Initially, floor-plan generation outlines the architectural skeleton, determining room configurations and the layout's overall flow. Subsequently, doorway and window placement modules add granularity, focusing on optimizing accessibility and lighting. Finally, the object selection and placement modules populate the environment, ensuring semantic consistency between objects and adhering to spatial constraints. 
